\documentclass[11pt,a4paper]{article}

\usepackage[margin=1in]{geometry}
\usepackage{amsmath,amssymb,amsthm}
\usepackage{bm}
\usepackage{booktabs}
\usepackage{hyperref}
\usepackage{enumitem}
\usepackage{xcolor}

\newcommand{\R}{\mathbb{R}}
\newcommand{\E}{\mathbb{E}}
\newcommand{\tr}{\operatorname{tr}}
\newcommand{\diag}{\operatorname{diag}}
\newcommand{\argmin}{\operatorname{argmin}}
\newcommand{\argmax}{\operatorname{argmax}}
\newcommand{\Cov}{\operatorname{Cov}}
\newcommand{\Var}{\operatorname{Var}}
\newcommand{\KL}{\operatorname{KL}}

\newtheorem{theorem}{Theorem}
\newtheorem{corollary}{Corollary}
\newtheorem{remark}{Remark}
\newtheorem{definition}{Definition}

\title{Adaptive Experimental Design:\\Pointwise Information, the Ignorance Gap,\\and Joint Geometry--Policy Optimization}
\author{Z.\ Lin}
\date{April 2026}

\begin{document}
\maketitle

\begin{abstract}
We develop a self-contained theoretical framework for the joint optimization of sensor hardware $\bm{c}$ and adaptive measurement policy $\pi$ under a Bayesian prior on the unknown state $\bm{x}$.  The framework is organized around the \emph{pointwise expected information gain} $\Phi(\bm{x}, \bm{s}, \bm{c})$---the expected reduction in posterior entropy about $\bm{x}$ under observations drawn from the true-$\bm{x}$ likelihood.  Its prior average is the standard mutual information $I(\bm{x};\,\bm{y}\mid\bm{s},\bm{c})$, and its pointwise evaluation cleanly defines the clairvoyant \emph{oracle} $V_{\mathrm{oracle}}(\bm{x}) = \max_{\bm{s}} \Phi(\bm{x}, \bm{s}, \bm{c})$, the pre-committed \emph{fixed design} $V_{\mathrm{fixed}}(\bm{c}) = \max_{\bm{s}} \E_{\bm{x}}[\Phi(\bm{x}, \bm{s}, \bm{c})]$, the \emph{ignorance gap} $IG(\bm{x}) = V_{\mathrm{oracle}}(\bm{x}) - \Phi(\bm{x},\bm{s}^\star_{\mathrm{fixed}},\bm{c})$, and the adaptive value $V_{\mathrm{adaptive}}(\bm{c})$ via Bellman dynamic programming on the belief state.  The prior-expected ignorance gap $\E[IG]$ upper-bounds the adaptive advantage $V_{\mathrm{adaptive}} - V_{\mathrm{fixed}}$---the classical Expected Value of Perfect Information (EVPI).  We show this bound is typically loose and, crucially, that $\argmax_{\bm{c}} \E[IG](\bm{c}) \neq \argmax_{\bm{c}} V_{\mathrm{adaptive}}(\bm{c})$: EVPI is a go/no-go diagnostic, not a design target.  We formulate the joint $(\bm{c}, \pi)$ optimization, derive the envelope-theorem geometry gradient that avoids differentiating through the policy's $\argmax$, and prescribe a policy-iteration scheme that alternates Bellman policy solves with gradient steps on $\bm{c}$.  A fully-solved radar beam-search case study exhibits every ingredient at minimum dimension, including a quantitative $\sim\!2.8\times$ separation between EVPI-guided and joint-DP-guided geometry selection.  A final section catalogs principled relaxations when the exact Bellman DP is intractable, ordered by structure preserved.  An appendix derives the Bellman equation from first principles.
\end{abstract}

\tableofcontents

%===============================================================================
\section{Introduction}
\label{sec:intro}
%===============================================================================

A sensor designed in isolation and a measurement policy designed in isolation are both half-designed.  When the operator of the sensor will react to intermediate observations---steering an antenna, changing a pulse delay, reconfiguring a metasurface---the optimal hardware depends on what the operator will actually do, and the optimal operator strategy depends on what the hardware allows.  This paper develops a clean formulation of that joint problem.

The organizing quantity throughout is not Fisher information, not the Cram\'er--Rao bound, not the posterior variance, but a single information-theoretic functional: the \emph{pointwise expected information gain} $\Phi(\bm{x}, \bm{s}, \bm{c})$, defined as the expected reduction in posterior entropy about the unknown state $\bm{x}$ when the measurement is performed with sensor parameters $\bm{s}$ and device geometry $\bm{c}$, \emph{given} that the true state is $\bm{x}$.  Its prior average is the standard mutual information between $\bm{x}$ and the data.  Its pointwise evaluation---holding $\bm{x}$ fixed---lets us define a clairvoyant oracle that knows $\bm{x}$, a non-adaptive fixed design that must hedge across the prior, and an adaptive policy that reacts to intermediate data.  The hierarchy
\[
  V_{\mathrm{fixed}}(\bm{c}) \;\le\; V_{\mathrm{adaptive}}(\bm{c}) \;\le\; \E_{\bm{x}}[V_{\mathrm{oracle}}(\bm{x};\,\bm{c})]
\]
is exact.  The gap between the first two is the quantity adaptation can win; the gap between the first and third is the ignorance-gap / EVPI upper bound on that win.

Two observations drive the rest of the paper.  First, the EVPI bound is often loose: typical geometries realize only a fraction of the bound as actual adaptive advantage, and some realize essentially none.  Second, the geometry $\bm{c}$ that maximizes $\E[IG](\bm{c})$ is generically \emph{not} the geometry that maximizes $V_{\mathrm{adaptive}}(\bm{c})$.  In our radar case study the two objectives pick geometries whose $V_{\mathrm{adaptive}}$ values differ by a factor of $\sim\!2.8$.  EVPI is a legitimate pre-solver diagnostic---if $\E[IG]$ is uniformly small across candidate $\bm{c}$, no adaptive policy can pay for itself---but it is not the right function to optimize $\bm{c}$ against.  The right function is $V_{\mathrm{adaptive}}(\bm{c})$, and it can only be computed by solving the adaptive problem at each candidate $\bm{c}$.

\paragraph{Structure of the paper.}
Section~\ref{sec:reward} introduces $\Phi(\bm{x}, \bm{s}, \bm{c})$ and justifies its use as the reward.  Sections~\ref{sec:fixed}--\ref{sec:ig} define the fixed design, the oracle, and the ignorance gap, and state the EVPI bound.  Sections~\ref{sec:adaptive}--\ref{sec:bellman} define adaptive policies and characterize the optimal policy via Bellman's equation on the belief state.  Section~\ref{sec:joint} gives the joint $(\bm{c}, \pi)$ formulation and its envelope-theorem solution.  Section~\ref{sec:examples} is the radar beam-search worked example.  Section~\ref{sec:relaxations} catalogs principled relaxations when exact Bellman is infeasible.  Appendix~\ref{app:bellman} derives the Bellman equation from first principles.

%===============================================================================
\section{The reward function: pointwise expected information gain}
\label{sec:reward}
%===============================================================================

\subsection{Setup}

Let $\bm{x} \in \mathcal{X}$ be an unknown state drawn from a prior $p_0(\bm{x})$.  A sensor produces observations $\bm{y}$ via a likelihood $p(\bm{y} \mid \bm{x}, \bm{s}, \bm{c})$, where $\bm{s} \in \R^{d_s}$ is the (possibly multi-step) sensor control and $\bm{c} \in \R^{d_c}$ is the device geometry held fixed across one sensing episode.  The posterior after observing $\bm{y}$ under $(\bm{s}, \bm{c})$ is
\begin{equation}
  p(\bm{x} \mid \bm{y}, \bm{s}, \bm{c}) \;=\; \frac{p(\bm{y} \mid \bm{x}, \bm{s}, \bm{c})\, p_0(\bm{x})}{p(\bm{y} \mid \bm{s}, \bm{c})},
  \qquad
  p(\bm{y} \mid \bm{s}, \bm{c}) \;=\; \int p(\bm{y} \mid \bm{x}', \bm{s}, \bm{c})\, p_0(\bm{x}')\, d\bm{x}'.
\end{equation}
Write $H(p) = -\int p(\bm{x}) \ln p(\bm{x})\, d\bm{x}$ for the differential entropy of a density $p$ (nats).

\subsection{Definition}

\begin{definition}[Pointwise expected information gain]
\label{def:phi}
For a fixed $\bm{x}\in\mathcal{X}$, the \emph{pointwise expected information gain} at $\bm{x}$ under sensor control $\bm{s}$ and geometry $\bm{c}$ is
\begin{equation}
  \boxed{\;
  \Phi(\bm{x}, \bm{s}, \bm{c}) \;=\; H(p_0) \;-\; \E_{\bm{y} \sim p(\bm{y}\mid \bm{x}, \bm{s}, \bm{c})}\!\bigl[\, H\!\bigl(p(\,\cdot\mid\bm{y}, \bm{s}, \bm{c})\bigr) \,\bigr].
  \;}
  \label{eq:phi_def}
\end{equation}
\end{definition}

In words: $\Phi(\bm{x}, \bm{s}, \bm{c})$ is the expected decrease in posterior entropy about $\bm{x}$ when the data are generated under the true-$\bm{x}$ model.  It conditions on $\bm{x}$---it is the information gain that results from an observation generated at $\bm{x}$, averaged over the randomness of the observation but not over $\bm{x}$ itself.

Equivalently (and sometimes more useful computationally), $\Phi(\bm{x}, \bm{s}, \bm{c})$ is the expected KL divergence from the posterior to the prior under the true-$\bm{x}$ data-generating distribution:
\begin{equation}
  \Phi(\bm{x}, \bm{s}, \bm{c}) \;=\; \E_{\bm{y} \sim p(\bm{y}\mid \bm{x}, \bm{s}, \bm{c})}\!\bigl[\,\KL\!\bigl(p(\,\cdot\mid\bm{y}, \bm{s}, \bm{c}) \,\big\|\, p_0\bigr)\,\bigr].
  \label{eq:phi_kl}
\end{equation}
The equivalence of \eqref{eq:phi_def} and \eqref{eq:phi_kl} follows from $\KL(p \,\|\, p_0) = -H(p) + H(p_0) + (\text{cross-entropy terms that integrate to zero})$; see Appendix~\ref{app:phi_identities}.

\subsection{Relation to standard mutual information}

Standard mutual information between $\bm{x}$ and $\bm{y}$ under $(\bm{s}, \bm{c})$ is a single scalar obtained by averaging over the joint distribution:
\begin{equation}
  I(\bm{x}; \bm{y} \mid \bm{s}, \bm{c})
  \;=\; H(p_0) - \E_{\bm{y} \sim p(\cdot \mid \bm{s}, \bm{c})}\!\bigl[H\!\bigl(p(\,\cdot\mid\bm{y},\bm{s},\bm{c})\bigr)\bigr]
  \;=\; \E_{\bm{x}, \bm{y}}\!\left[\ln \frac{p(\bm{y}\mid\bm{x},\bm{s},\bm{c})}{p(\bm{y}\mid\bm{s},\bm{c})}\right].
\end{equation}
Mutual information is not a function of $\bm{x}$.

\begin{theorem}[Prior average of $\Phi$ equals MI]
\label{thm:phi_is_mi}
\begin{equation}
  \E_{\bm{x} \sim p_0}\!\bigl[\Phi(\bm{x}, \bm{s}, \bm{c})\bigr] \;=\; I(\bm{x}; \bm{y} \mid \bm{s}, \bm{c}).
  \label{eq:phi_average}
\end{equation}
\end{theorem}

\begin{proof}
Apply $\E_{\bm{x}}$ to \eqref{eq:phi_def} and use $\E_{\bm{x}} \E_{\bm{y}\mid\bm{x}}[f(\bm{y})] = \E_{\bm{y}}[f(\bm{y})]$ where the last expectation is under the marginal $p(\bm{y}\mid\bm{s},\bm{c})$.  The prior entropy $H(p_0)$ is constant.\end{proof}

So $\Phi(\bm{x}, \bm{s}, \bm{c})$ is exactly the \emph{pointwise building block whose prior-average is mutual information}.  Two uses in the rest of the paper:

\begin{itemize}[nosep, leftmargin=2em]
  \item Oracle / ignorance gap analyses hold $\bm{x}$ fixed; the right quantity there is $\Phi(\bm{x}, \cdot, \cdot)$ (Sections~\ref{sec:ig}).
  \item Fixed and adaptive designs integrate over $\bm{x}$; the right quantity there is $\E_{\bm{x}}[\Phi(\bm{x}, \cdot, \cdot)]$, which by Theorem~\ref{thm:phi_is_mi} is MI (Sections~\ref{sec:fixed}, \ref{sec:adaptive}--\ref{sec:bellman}).
\end{itemize}

\subsection{Non-negativity and basic properties}

\begin{remark}[Non-negativity in expectation]
\label{rem:nonneg}
By Jensen's inequality and $\E_{\bm{y}\mid\bm{x}}[H(p(\,\cdot\mid\bm{y}))] = H(p_0) - \Phi(\bm{x},\bm{s},\bm{c})$, one has $\E_{\bm{x}}[\Phi] = I \ge 0$.  Pointwise $\Phi(\bm{x},\bm{s},\bm{c})$ can in principle be negative for a specific $\bm{x}$ (the data generated at an ``unusual'' $\bm{x}$ may, in expectation, produce a posterior that is more diffuse than the prior in entropy terms), but the prior-weighted average is non-negative.  In the regular, well-identifiable models used throughout this paper, pointwise $\Phi(\bm{x},\bm{s},\bm{c}) \ge 0$ for all $\bm{x}$; negative values almost always indicate a bug or a genuinely pathological likelihood.
\end{remark}

\begin{remark}[Multi-step $\Phi$]
\label{rem:multistep}
For $N$ conditionally independent measurements under a sequence $\bm{s}_{1:N}$,
\begin{equation}
  \Phi(\bm{x}, \bm{s}_{1:N}, \bm{c}) \;=\; H(p_0) \;-\; \E_{\bm{y}_{1:N}\mid \bm{x}, \bm{s}_{1:N}, \bm{c}}\!\bigl[\, H\!\bigl(p(\,\cdot\mid\bm{y}_{1:N},\bm{s}_{1:N},\bm{c})\bigr)\,\bigr],
\end{equation}
with the posterior updated by the full data $\bm{y}_{1:N}$.  All arguments below are phrased for the $N$-step reward.
\end{remark}

\subsection{Gaussian-posterior / Fisher specialization}
\label{subsec:fisher_specialization}

For readers coming from classical experimental-design literature, it is useful to note the specialization of $\Phi$ in the Gaussian-posterior / Bernstein--von Mises regime.  Suppose the observations are sufficiently informative that, conditional on the true $\bm{x}_0$, the posterior $p(\,\cdot\mid\bm{y}_{1:N})$ concentrates around $\bm{x}_0$ and is approximately Gaussian with covariance $\bm{J}_N(\bm{x}_0, \bm{s}_{1:N}, \bm{c})^{-1}$, where $\bm{J}_N$ is the usual Fisher information matrix accumulated over the $N$ measurements.  Under this approximation, $H(\text{posterior}) \approx \tfrac{1}{2} \ln\det(2\pi e \bm{J}_N^{-1})$, so
\begin{equation}
  \Phi(\bm{x}, \bm{s}_{1:N}, \bm{c}) \;\approx\; \mathrm{const} \;+\; \tfrac{1}{2} \ln\det \bm{J}_N(\bm{x}, \bm{s}_{1:N}, \bm{c}).
  \label{eq:phi_fisher}
\end{equation}
That is, in the Gaussian / high-SNR regime, $\Phi$ specializes to the log-volume of the Fisher information ellipsoid---the classical \emph{D-optimal} criterion of experimental design.  All subsequent developments in this paper (oracle, ignorance gap, Bellman DP, joint optimization) are phrased in terms of $\Phi$ and therefore include \eqref{eq:phi_fisher} as a special case.  Outside the Gaussian regime---multi-modal or heavy-tailed posteriors, strong prior regularization---one must use $\Phi$ directly rather than its Fisher surrogate; the approximation \eqref{eq:phi_fisher} breaks down.

Throughout the main text we work with $\Phi$ directly.  When a case study permits the Gaussian specialization (Section~\ref{sec:relaxations}, path (2c)), we flag it explicitly.

\subsection{Why $\Phi$ is the right reward}

\begin{itemize}[nosep, leftmargin=2em]
  \item \textbf{Information content is what a sensor produces.}  $\Phi$ is a property of the joint distribution $p(\bm{x}, \bm{y} \mid \bm{s}, \bm{c})$ induced by the sensor model.  No estimator is involved in defining or computing it.  It is the quantity that determines---in aggregate and through the posterior---what any downstream estimator can achieve.
  \item \textbf{It supports the oracle / ignorance-gap decomposition.}  The pointwise form is essential: $V_{\mathrm{oracle}}(\bm{x})$ and $IG(\bm{x})$ only make sense as functions of $\bm{x}$.  Standard MI, which is a single scalar, cannot play this role.
  \item \textbf{It reduces to classical D-optimality in the Gaussian/BvM regime.}  Equation~\eqref{eq:phi_fisher} shows the framework is backward-compatible with decades of experimental-design literature, without being dependent on it.
  \item \textbf{Estimator-agnostic.}  Debates about Cram\'er--Rao bounds, posterior covariance, James--Stein-type shrinkage, and biased-vs-unbiased estimators are all downstream of $\Phi$.  Maximizing $\Phi$ concentrates the posterior; how a particular estimator extracts a point from that posterior is an orthogonal question.
\end{itemize}

%===============================================================================
\section{Why information gain and not PCRB? A structural asymmetry}
\label{sec:phi_vs_pcrb}
%===============================================================================

A natural alternative to the framework just laid out would build on the prior-averaged posterior Cram\'er--Rao bound (PCRB) rather than on the pointwise information gain $\Phi$.  Both formulations admit a clean pre-committed design problem.  The rest of the machinery---oracle, ignorance gap, Bellman reward at a realized trajectory---looks similar in structure.  This section shows that the similarity is superficial: only the $\Phi$ framework admits a rigorous pointwise extension.  The PCRB-based alternative is strictly weaker, by two independent mechanisms.  The conclusion is not that PCRB is wrong (it is a correct and useful bound), but that it is insufficient as the foundational object of a joint geometry--policy theory.

\subsection{The PCRB-based fixed-design alternative}
\label{subsec:pcrb_alternative}

For definiteness, consider a static unknown state $\bm{x} \sim p_0$ observed through $N$ conditionally independent measurements $\bm{y}_k = \bm{h}(\bm{x}, \bm{s}_k, \bm{c}) + \bm{v}_k$, $\bm{v}_k \sim \mathcal{N}(\bm{0}, \sigma^2 \bm{I})$, with Jacobian $\bm{F}_k(\bm{x}) = \partial\bm{h}/\partial\bm{x}$.  Accumulating Fisher information across measurements and averaging over the prior produces the Bayesian Information Matrix
\begin{equation}
  \bm{J}_P(\bm{s}_{1:N}, \bm{c}) \;=\; \bm{J}_0 \;+\; \sigma^{-2}\sum_{k=1}^N \E_{\bm{x} \sim p_0}\!\bigl[\bm{F}_k(\bm{x}, \bm{s}_k, \bm{c})^\top \bm{F}_k(\bm{x}, \bm{s}_k, \bm{c})\bigr],
  \label{eq:JP}
\end{equation}
where $\bm{J}_0 = -\E_{\bm{x}}[\nabla^2 \ln p_0(\bm{x})]$ is the prior information.  The PCRB (Van Trees) states that for any estimator $\hat{\bm{x}}(\bm{y}_{1:N})$,
\begin{equation}
  \E_{\bm{x}, \bm{y}_{1:N}}\!\bigl[(\hat{\bm{x}} - \bm{x})(\hat{\bm{x}} - \bm{x})^\top\bigr] \;\succeq\; \bm{J}_P(\bm{s}_{1:N}, \bm{c})^{-1}.
  \label{eq:pcrb_ineq}
\end{equation}
Minimizing the log-volume of the PCRB ellipsoid gives the Bayesian D-optimal fixed-design problem
\begin{equation}
  V^{\mathrm{PCRB}}_{\mathrm{fixed}}(\bm{c}) \;=\; \max_{\bm{s}_{1:N}}\; \ln\det \bm{J}_P(\bm{s}_{1:N}, \bm{c}).
  \label{eq:vfixed_pcrb}
\end{equation}
This is a clean single-level, estimator-free objective---structurally parallel to the $\Phi$-based $V_{\mathrm{fixed}}(\bm{c}) = \max_{\bm{s}} \E_{\bm{x}}[\Phi(\bm{x}, \bm{s}, \bm{c})]$ of the next section.

If the framework stopped at the fixed-design problem, the two would be roughly interchangeable (the $\Phi$ formulation reduces to \eqref{eq:vfixed_pcrb} in the Gaussian-posterior regime per \eqref{eq:phi_fisher}).  The divergence begins when one attempts to extend the framework to a pointwise oracle, a pointwise ignorance gap, or a Bellman terminal reward.

\subsection{Pointwise values can beat the aggregate in both frameworks}
\label{subsec:pointwise_both}

A tempting way to distinguish the two frameworks is to observe that pointwise values at specific $\bm{x}$ can exceed the global prior-averaged bound.  This is \emph{not} a distinguishing feature---it happens in both frameworks:

\begin{itemize}[nosep, leftmargin=2em]
  \item In the $\Phi$ framework, $\Phi(\bm{x}, \bm{s}, \bm{c})$ is the expected information gain \emph{conditional on $\bm{x}$}.  By Theorem~\ref{thm:phi_is_mi}, $\E_{\bm{x}}[\Phi(\bm{x},\bm{s},\bm{c})] = I(\bm{x};\bm{y}\mid\bm{s},\bm{c})$.  Pointwise $\Phi(\bm{x})$ can exceed or fall short of $I$; the inequality holds only in the prior-weighted integral.
  \item In the PCRB framework, the pointwise MSE of a Bayesian estimator (posterior mean, MAP) at a specific $\bm{x}_0$ can be \emph{below} $\tr \bm{J}_P^{-1}$, at the cost of being above $\tr \bm{J}_P^{-1}$ at other $\bm{x}$ values.  The prior acts as shrinkage: it reduces variance at favorable $\bm{x}$ (e.g., interior of the prior support) while biasing estimates at unfavorable $\bm{x}$ (e.g., boundary).  The PCRB inequality \eqref{eq:pcrb_ineq} holds only for the prior-weighted integral.
\end{itemize}

The question is not whether pointwise values can beat aggregate bounds (yes in both), but whether each framework supplies a \emph{rigorous pointwise object} that can serve as the building block for an oracle or ignorance-gap analysis.

\subsection{The asymmetry: \texorpdfstring{$\Phi$}{Phi} has a pointwise summand, PCRB does not}
\label{subsec:asymmetry}

\paragraph{In the $\Phi$ framework, $\Phi(\bm{x})$ is literally the pointwise integrand of MI.}  By Theorem~\ref{thm:phi_is_mi},
\begin{equation}
  I(\bm{x};\bm{y}\mid\bm{s},\bm{c}) \;=\; \int \Phi(\bm{x}, \bm{s}, \bm{c})\, p_0(\bm{x})\, d\bm{x}.
\end{equation}
$\Phi(\bm{x})$ is the rigorous pointwise object whose prior-weighted integral is the aggregate.  The construction is exact at the level of joint distributions $p(\bm{x}, \bm{y} \mid \bm{s}, \bm{c})$: no estimator is invoked, no approximation is made, and $\Phi(\bm{x})$ is well-defined for arbitrary (possibly highly non-Gaussian) likelihoods.  It supports unambiguously the oracle $V_{\mathrm{oracle}}(\bm{x};\bm{c}) = \max_{\bm{s}} \Phi(\bm{x}, \bm{s}, \bm{c})$, the pointwise ignorance gap $IG(\bm{x};\bm{c})$, and the Bellman terminal reward $R(b_{N+1}) = H(p_0) - H(b_{N+1})$.

\paragraph{The PCRB framework has no analogous pointwise summand.}  Two independent obstructions preclude a direct pointwise version:

\begin{enumerate}[leftmargin=2em, nosep]

\item \textbf{PCRB itself has no pointwise integrand.}  The bound \eqref{eq:pcrb_ineq} is formed by averaging Fisher information over the prior \emph{before} inversion: $\bm{J}_P = \bm{J}_0 + \sum_k \E_{\bm{x}}[\bm{F}_k^\top \bm{F}_k]$.  The natural candidate for a pointwise summand would be $\tr \bm{J}_N(\bm{x})^{-1}$, the inverse Fisher information at a specific $\bm{x}$.  But
\begin{equation}
  \E_{\bm{x}}\!\bigl[\tr \bm{J}_N(\bm{x})^{-1}\bigr] \;\neq\; \tr \bm{J}_P^{-1}
  \qquad\text{in general,}
\end{equation}
by Jensen's inequality (the map $\bm{A} \mapsto \tr \bm{A}^{-1}$ is convex on the positive-definite cone).  There is no decomposition of the PCRB bound into a prior-weighted integral of pointwise bounds.  The object $\tr \bm{J}_N(\bm{x})^{-1}$, if used at all, is an \emph{ad hoc} choice that does not recover \eqref{eq:pcrb_ineq} upon averaging.

\item \textbf{No pointwise MSE bound holds for biased Bayesian estimators.}  Even setting aside (1) and hoping that $\tr \bm{J}_N(\bm{x}_0)^{-1}$ bounds pointwise MSE at $\bm{x}_0$, one runs into the classical CRB obstruction.  The pointwise classical CRB
\begin{equation}
  \Cov_{\bm{y}\mid\bm{x}_0}\!\bigl[\hat{\bm{x}}(\bm{y})\bigr] \;\succeq\; \bm{J}_N(\bm{x}_0)^{-1}
  \label{eq:classical_crb}
\end{equation}
holds \emph{only for unbiased estimators} of the deterministic parameter $\bm{x}_0$.  Sequential Bayesian estimators---posterior mean, MAP, any Kalman-filter-like recursive estimator---are \emph{not} unbiased: the prior regularizes them, inducing a shrinkage bias $\bm{b}(\bm{x}_0) = \E[\hat{\bm{x}}\mid \bm{x}_0] - \bm{x}_0 \neq \bm{0}$ that is the price of admitting the prior at all.  Therefore \eqref{eq:classical_crb} does not apply.

The biased-CRB correction does exist,
\begin{equation}
  \Cov_{\bm{y}\mid\bm{x}_0}\!\bigl[\hat{\bm{x}}(\bm{y})\bigr] \;\succeq\; \bigl(\bm{I} + \nabla_{\bm{x}}\bm{b}(\bm{x}_0)\bigr)\, \bm{J}_N(\bm{x}_0)^{-1}\, \bigl(\bm{I} + \nabla_{\bm{x}}\bm{b}(\bm{x}_0)\bigr)^\top,
\end{equation}
but requires the gradient $\nabla_{\bm{x}}\bm{b}$ of the estimator's bias function---a quantity available in closed form only in trivial cases.  In any practical sequential-sensing problem the bias function is implicit, defined operationally by the recursive posterior-mean update, and its gradient is neither tabulated nor tractable.  The biased-CRB is therefore of theoretical but not operational value.

\end{enumerate}

\paragraph{Together.}  Obstruction (1) says PCRB does not provide a pointwise summand; obstruction (2) says no other pointwise MSE bound applies to the estimators that sequential Bayesian sensing actually uses.  The gap cannot be closed by more careful accounting.

\subsection{Asymptotic recovery is not enough}
\label{subsec:asymptotic}

In the Bernstein--von Mises regime (high SNR, large $N$, smooth identifiable model), the posterior concentrates near the true $\bm{x}_0$ with covariance approximately $\bm{J}_N(\bm{x}_0)^{-1}$, and the Bayesian estimator's pointwise MSE is $\bm{J}_N(\bm{x}_0)^{-1}$ \emph{asymptotically}.  This is an equality in the limit, not a bound at finite $N$.

For the joint-DP framework we need a rigorous pointwise object at \emph{all} $N$---oracle and ignorance-gap analyses operate precisely in the regime where the sensor budget is finite and the posterior still carries meaningful uncertainty that the policy must navigate.  Asymptotic equality at infinite $N$ does not supply the decomposition we need at the finite horizons of Sections~\ref{sec:bellman}--\ref{sec:examples}.  $\Phi$ does, at every $N$, without any asymptotic assumption.

\subsection{What PCRB can and cannot support}
\label{subsec:pcrb_capabilities}

The asymmetry has concrete consequences for what each framework can rigorously formulate.

\paragraph{PCRB rigorously supports:}
\begin{itemize}[nosep, leftmargin=2em]
  \item \textbf{Prior-averaged fixed design.}  $V^{\mathrm{PCRB}}_{\mathrm{fixed}}(\bm{c}) = \max_{\bm{s}} \ln\det \bm{J}_P$ is a well-posed single-level problem.
  \item \textbf{Gradient-based design.}  $\partial V^{\mathrm{PCRB}}_{\mathrm{fixed}} / \partial\bm{c}$ and $\partial V^{\mathrm{PCRB}}_{\mathrm{fixed}} / \partial\bm{s}_k$ are analytic in closed form.
  \item \textbf{Prior-averaged monotone bounds.}  Inequalities like $\E[\mathrm{MSE}_{\mathrm{fixed}}] \ge \tr\bm{J}_P^{-1}$.
\end{itemize}

\paragraph{PCRB cannot rigorously support:}
\begin{itemize}[nosep, leftmargin=2em]
  \item \textbf{Pointwise oracle $V_{\mathrm{oracle}}(\bm{x})$.}  Would require a pointwise MSE bound; obstructions (1) and (2) above preclude it.
  \item \textbf{Pointwise ignorance gap $IG(\bm{x})$.}  Defined as $V_{\mathrm{oracle}}(\bm{x}) - (\text{fixed value at } \bm{x})$; inherits the obstruction.
  \item \textbf{Bellman terminal reward at a realized trajectory.}  The terminal $\E_{\bm{x}\mid b_{N+1}}[\Phi(\bm{x}, \bm{s}_{1:N})]$ in \eqref{eq:terminal} requires pointwise evaluation of the reward along each observation path.  PCRB provides no such pointwise quantity.
\end{itemize}

These are not omissions that additional work could fill in; they are structurally absent because the PCRB inequality is an \emph{integral} bound by construction, with no pointwise decomposition surviving the prior-averaging step that defines $\bm{J}_P$.

\subsection{The \texorpdfstring{$\Phi$}{Phi} framework subsumes the PCRB framework in the Gaussian regime}
\label{subsec:phi_subsumes}

In the Gaussian-posterior / BvM regime of Section~\ref{subsec:fisher_specialization},
\begin{equation}
  \Phi(\bm{x}, \bm{s}, \bm{c}) \;\approx\; \mathrm{const} \;+\; \tfrac{1}{2}\,\ln\det \bm{J}_N(\bm{x}, \bm{s}, \bm{c}).
\end{equation}
Taking prior averages and applying the $\Phi$-to-MI identity:
\begin{equation}
  V_{\mathrm{fixed}}(\bm{c}) \;=\; \max_{\bm{s}}\; \E_{\bm{x}}[\Phi(\bm{x}, \bm{s}, \bm{c})] \;\approx\; \max_{\bm{s}}\; \tfrac{1}{2}\,\E_{\bm{x}}[\ln\det \bm{J}_N(\bm{x}, \bm{s}, \bm{c})] \;+\; \mathrm{const}.
  \label{eq:vfixed_gaussian}
\end{equation}
This is closely related to \eqref{eq:vfixed_pcrb}, the PCRB-based objective, though the two differ by the order of log-det and expectation (Jensen gap).  Both recover the classical Bayesian D-optimal program in the Gaussian regime, up to this reordering.

The $\Phi$ framework therefore includes the PCRB-style fixed-design objective as a special case (via \eqref{eq:vfixed_gaussian}), extends it pointwise (via $\Phi(\bm{x}, \cdot)$), and remains rigorous outside the Gaussian regime (exact MI without any approximation).  The reverse containment does not hold: PCRB admits no pointwise extension, and ceases to be a bound outside the regime where the estimators happen to be asymptotically unbiased.

\subsection{PCRB as an independent MSE-anchored baseline}
\label{subsec:pcrb_baseline}

The structural limitations of Section~\ref{subsec:asymmetry} do \emph{not} disqualify the PCRB framework from playing a useful role in the joint-design picture---they only prevent it from serving as the \emph{foundation}.  One constructive use remains: the PCRB-based joint fixed-design optimization \eqref{eq:vfixed_pcrb} serves as an \textbf{independent, MSE-anchored reference baseline} against which the $\Phi$-based joint $(\bm{c}, \pi)$ DP result can be reported.

The baseline is well-defined without any of the obstructions discussed above:
\begin{itemize}[nosep, leftmargin=2em]
  \item The PCRB inequality \eqref{eq:pcrb_ineq} is rigorous at any finite $N$ for \emph{any} estimator satisfying the mild regularity conditions of the bound.  No unbiasedness is assumed---the inequality is in prior-averaged form, which tolerates the bias inherent in Bayesian estimators because the expectation smooths over $\bm{x}$.
  \item The joint optimization
  \begin{equation}
    \bigl(\bm{c}_{\mathrm{PCRB}}^\star,\, \bm{s}^\star_{\mathrm{PCRB}}\bigr) \;=\; \argmax_{\bm{c},\, \bm{s}_{1:N}}\; \ln\det \bm{J}_P(\bm{s}_{1:N}, \bm{c})
    \label{eq:pcrb_baseline}
  \end{equation}
  is a single-level problem in $d_c + N d_s$ real variables---flat, amenable to gradient descent, and well-studied in the classical experimental-design literature.  It produces a geometry $\bm{c}_{\mathrm{PCRB}}^\star$ that rigorously minimizes the log-volume of the prior-averaged MSE ellipsoid.
  \item The output geometry carries a concrete MSE interpretation: any estimator deployed at $\bm{c}_{\mathrm{PCRB}}^\star$ with fixed schedule $\bm{s}^\star_{\mathrm{PCRB}}$ has prior-averaged MSE bounded below by $\tr \bm{J}_P(\bm{s}^\star_{\mathrm{PCRB}}, \bm{c}_{\mathrm{PCRB}}^\star)^{-1}$.
\end{itemize}

This baseline has three operational uses in reporting joint-$(\bm{c}, \pi)$ DP results.

\paragraph{Sanity check on adaptive gain.}
Evaluating $V_{\mathrm{adaptive}}$ at $\bm{c}_{\mathrm{PCRB}}^\star$ (the baseline geometry) isolates the policy contribution: if the adaptive policy at the PCRB-optimal geometry already substantially beats the fixed $\bm{s}^\star_{\mathrm{PCRB}}$, we learn that adaptation pays off even on a geometry designed for the non-adaptive setting.  If not, the case for joint-policy optimization is weak at this geometry, independent of whether a better $\bm{c}$ exists.

\paragraph{MSE-anchored quantitative claim.}
The comparison
\begin{equation}
  \text{MSE ratio} \;=\; \frac{\tr \bm{\Sigma}_N^{\mathrm{adaptive}}(\bm{c}^\star)}{\tr \bm{J}_P(\bm{s}^\star_{\mathrm{PCRB}}, \bm{c}_{\mathrm{PCRB}}^\star)^{-1}}
\end{equation}
gives a reviewer-friendly quantitative statement of the combined joint-geometry-and-adaptive-policy gain, in concrete prior-averaged MSE units.  Unlike an information-gain ratio, which invites the question ``but what does this buy me in MSE?'', the PCRB-anchored ratio answers that question directly.

\paragraph{Null hypothesis for the whole framework.}
If $V_{\mathrm{adaptive}}(\bm{c}^\star)$ does not exceed $V^{\mathrm{PCRB}}_{\mathrm{fixed}}(\bm{c}_{\mathrm{PCRB}}^\star)$ after the Gaussian-regime translation of \eqref{eq:vfixed_gaussian}, then the joint-$(\bm{c}, \pi)$ DP machinery does not pay for itself relative to classical Bayesian D-optimal design.  In that case, one should report the PCRB-optimal $\bm{c}_{\mathrm{PCRB}}^\star$ with its fixed schedule and stop.  This is the quantitative analog of Corollary~\ref{cor:useless_adapt}, operationalized in MSE units.

\paragraph{Summary of the relationship.}
The two frameworks are \emph{complementary}, with distinct scopes:

\begin{itemize}[nosep, leftmargin=2em]
  \item The $\Phi$ framework supplies the \emph{internal machinery}---oracle, ignorance gap, Bellman reward, envelope-theorem gradient, joint optimization.  It is rigorous pointwise, estimator-agnostic, and valid at all $N$ and all likelihoods.
  \item The PCRB framework supplies the \emph{external MSE anchor}---a rigorously prior-averaged baseline geometry $\bm{c}_{\mathrm{PCRB}}^\star$ and a concrete MSE bound against which the $\Phi$-framework's output can be benchmarked.
\end{itemize}

Neither replaces the other.  The $\Phi$ framework without the PCRB baseline lacks a concrete MSE endpoint for its information-theoretic quantities; the PCRB framework without $\Phi$ lacks the pointwise machinery needed to formulate the adaptive half of the problem.  Reporting the joint-DP result and the PCRB baseline side-by-side is the natural way to close both gaps.

\subsection{Summary}
\label{subsec:phi_vs_pcrb_summary}

\begin{center}
\begin{tabular}{@{}lll@{}}
\toprule
Capability & $\Phi$ framework & PCRB framework \\
\midrule
Prior-averaged fixed design & yes, rigorous at any $N$ & yes, rigorous at any $N$ \\
Pointwise reward at specific $\bm{x}$ & $\Phi(\bm{x}, \bm{s}, \bm{c})$ (exact integrand of MI) & ad hoc; no rigorous candidate \\
Valid for biased estimators & yes (estimator-agnostic) & no; bound requires unbiasedness \\
Oracle $V_{\mathrm{oracle}}(\bm{x})$ & well-defined pointwise & not available \\
Ignorance gap $IG(\bm{x})$ & well-defined pointwise & not available \\
Bellman terminal reward at $b_{N+1}$ & $H(p_0) - H(b_{N+1})$ & no pointwise analog \\
Gaussian/BvM specialization & recovers Bayesian D-optimal & (definitional) \\
Validity outside Gaussian regime & exact & bound does not apply \\
\bottomrule
\end{tabular}
\end{center}

Three observations consolidate the asymmetry.  First, both frameworks admit the Jensen-gap phenomenon in which pointwise values can exceed the aggregate bound; this is neither special nor distinguishing.  Second, $\Phi(\bm{x})$ admits a rigorous pointwise interpretation as the pointwise integrand of MI, whereas no such pointwise integrand of PCRB exists.  Third, even if one wished to bound pointwise MSE directly, no such non-asymptotic bound applies to the biased Bayesian estimators actually used in sequential sensing.  The $\Phi$ framework is therefore the minimal rigorous setting in which the joint geometry--policy theory---oracle, ignorance gap, Bellman reward, envelope-theorem gradient, and joint optimization---can be stated at all.

%===============================================================================
\section{Pre-computed sensors: the fixed-design problem}
\label{sec:fixed}
%===============================================================================

A \emph{pre-computed} (or \emph{fixed}, or \emph{non-adaptive}) design commits the sensor sequence $\bm{s}_{1:N}$ at design time, before any data are collected.  Each $\bm{s}_k$ is a fixed vector; none depends on $\bm{y}_{1:k-1}$.

\begin{definition}[Fixed-design value]
\label{def:vfixed}
The fixed-design value at geometry $\bm{c}$ is
\begin{equation}
  V_{\mathrm{fixed}}(\bm{c}) \;=\; \max_{\bm{s}_{1:N}} \; \E_{\bm{x} \sim p_0}\bigl[\,\Phi(\bm{x}, \bm{s}_{1:N}, \bm{c})\,\bigr]
  \;=\; \max_{\bm{s}_{1:N}} \; I(\bm{x};\, \bm{y}_{1:N} \mid \bm{s}_{1:N}, \bm{c}),
  \label{eq:fixed}
\end{equation}
with optimizer $\bm{s}^\star_{\mathrm{fixed}}(\bm{c}) = \argmax_{\bm{s}_{1:N}} \E_{\bm{x}}[\Phi(\bm{x}, \bm{s}_{1:N}, \bm{c})]$.
\end{definition}

The equality with mutual information on the right is Theorem~\ref{thm:phi_is_mi}; the fixed-design problem is a single-level maximization of $I(\bm{x};\bm{y}_{1:N})$ over the $N d_s$-dimensional vector of sensor parameters (plus any geometry $\bm{c}$ one includes in the outer max).

\paragraph{Optimization structure.}
The fixed-design problem is a \emph{single-level} optimization---no inner argmax inside an outer max.  If $\Phi$ is smooth in $(\bm{s}, \bm{c})$, gradient-based optimization applies directly:
\begin{equation}
  \nabla_{\bm{s}_k} V_{\mathrm{fixed}}(\bm{c}) \;=\; \E_{\bm{x}}\!\bigl[\nabla_{\bm{s}_k} \Phi(\bm{x}, \bm{s}_{1:N}, \bm{c})\bigr]_{\bm{s}^\star},
  \qquad
  \nabla_{\bm{c}} V_{\mathrm{fixed}}(\bm{c}) \;=\; \E_{\bm{x}}\!\bigl[\nabla_{\bm{c}} \Phi(\bm{x}, \bm{s}_{1:N}, \bm{c})\bigr]_{\bm{s}^\star}.
\end{equation}
Monte Carlo or grid-based estimation of $\E_{\bm{x}}[\cdot]$ typically suffices.

\paragraph{Historical note.}
In the Gaussian-posterior specialization \eqref{eq:phi_fisher}, \eqref{eq:fixed} reduces to the classical Bayesian D-optimal design problem,
\begin{equation}
  V_{\mathrm{fixed}}^{\text{Fisher}}(\bm{c}) \;=\; \max_{\bm{s}_{1:N}} \; \E_{\bm{x}}\!\bigl[\,\tfrac{1}{2}\ln\det \bm{J}_N(\bm{x}, \bm{s}_{1:N}, \bm{c})\,\bigr] \;+\; \mathrm{const},
  \label{eq:joint_opt}
\end{equation}
recovering a century of sequential-design literature.  We keep \eqref{eq:fixed} as the primary statement because it is valid beyond the Gaussian regime.

%===============================================================================
\section{The ignorance gap}
\label{sec:ig}
%===============================================================================

The fixed design \eqref{eq:fixed} hedges against the unknown $\bm{x}$ by averaging over the prior.  How much worse is this hedging than an oracle that \emph{sees} $\bm{x}$ and picks sensors specifically for it?  The ignorance gap quantifies this pointwise.

\subsection{Oracle value and ignorance gap}

\begin{definition}[Oracle value]
\label{def:oracle}
The oracle value at state $\bm{x}$ is
\begin{equation}
  V_{\mathrm{oracle}}(\bm{x};\, \bm{c}) \;=\; \max_{\bm{s}_{1:N}} \; \Phi(\bm{x}, \bm{s}_{1:N}, \bm{c}).
  \label{eq:oracle}
\end{equation}
\end{definition}

This is the best expected information gain at $\bm{x}$ if the designer were allowed to see $\bm{x}$ and pick sensors specifically for it.  It is pointwise in $\bm{x}$.  No realizable policy (adaptive or otherwise) can exceed it, because no realizable policy has access to $\bm{x}$ before data collection.

\begin{definition}[Ignorance gap]
\label{def:ig}
The pointwise \emph{ignorance gap} at $\bm{x}$ is
\begin{equation}
  \boxed{\;
    IG(\bm{x}; \bm{c}) \;=\; V_{\mathrm{oracle}}(\bm{x}; \bm{c}) \;-\; \Phi\!\bigl(\bm{x},\, \bm{s}^\star_{\mathrm{fixed}}(\bm{c}),\, \bm{c}\bigr) \;\ge\; 0.
  \;}
  \label{eq:IG}
\end{equation}
\end{definition}

$IG(\bm{x})$ measures how much better the sensors \emph{could} be at $\bm{x}$ if tailored specifically to it, relative to the one-size-fits-all $\bm{s}^\star_{\mathrm{fixed}}$.  The prior-expected ignorance gap is
\begin{equation}
  \E_{\bm{x}}[IG(\bm{x}; \bm{c})] \;=\; \E_{\bm{x}}[V_{\mathrm{oracle}}(\bm{x}; \bm{c})] \;-\; V_{\mathrm{fixed}}(\bm{c}).
  \label{eq:evpi_identity}
\end{equation}

\subsection{The EVPI bound}

Introducing (a preview of Section~\ref{sec:adaptive}) the value of the optimal adaptive policy $V_{\mathrm{adaptive}}(\bm{c})$:

\begin{theorem}[EVPI bound]
\label{thm:evpi}
\begin{equation}
  \boxed{\;0 \;\le\; V_{\mathrm{adaptive}}(\bm{c}) - V_{\mathrm{fixed}}(\bm{c}) \;\le\; \E_{\bm{x}}[IG(\bm{x}; \bm{c})].\;}
  \label{eq:evpi}
\end{equation}
\end{theorem}

\begin{proof}
Lower bound: a fixed schedule is a (constant) adaptive policy, so $V_{\mathrm{adaptive}} \ge V_{\mathrm{fixed}}$.  Upper bound: no policy has access to $\bm{x}$, so $V_{\mathrm{adaptive}} \le \E_{\bm{x}}[V_{\mathrm{oracle}}(\bm{x})]$; subtract $V_{\mathrm{fixed}}$ from both sides and apply \eqref{eq:evpi_identity}.
\end{proof}

This is the \emph{Expected Value of Perfect Information} (EVPI) bound from decision theory, specialized to information-gain rewards.

\begin{corollary}
\label{cor:useless_adapt}
If $\E_{\bm{x}}[IG(\bm{x}; \bm{c})] \approx 0$, then no adaptive policy can significantly outperform the pre-computed design at that $\bm{c}$.
\end{corollary}

\begin{remark}[EVPI is necessary but not sufficient for adaptive gain]
\label{rem:evpi_necessary_not_sufficient}
The converse of Corollary~\ref{cor:useless_adapt} does \emph{not} hold.  A large $\E[IG]$ indicates that adaptation \emph{might} help but does not guarantee it.  The inequality \eqref{eq:evpi} is generically loose, and the realized gap $V_{\mathrm{adaptive}} - V_{\mathrm{fixed}}$ depends sensitively on the geometry $\bm{c}$.  Two designs with the same $\E[IG]$ can have radically different $V_{\mathrm{adaptive}} - V_{\mathrm{fixed}}$; a design with smaller $\E[IG]$ can have a \emph{larger} realized adaptive advantage than one with larger $\E[IG]$.  Section~\ref{subsec:value_of_joint} quantifies this in the radar case study, where the EVPI-maximizing geometry has $\sim\!2.8\times$ \emph{smaller} $V_{\mathrm{adaptive}}$ than the joint-DP optimum within the same beam family.  Consequently, $\E[IG]$ is a pre-solver diagnostic, not a design target: it says \emph{whether} to engage the adaptive loop, not \emph{where} the joint $(\bm{c}, \pi)$ optimum lies.
\end{remark}

\subsection{When is the ignorance gap small?}

$IG(\bm{x}) \equiv 0$ exactly when $\argmax_{\bm{s}_{1:N}} \Phi(\bm{x}, \bm{s}_{1:N}, \bm{c})$ does not depend on $\bm{x}$---i.e., the same sensor sequence is oracle-optimal regardless of the true state.  This holds for \emph{translation-covariant} likelihoods where $\bm{x}$ shifts the data but does not change which sensor configuration extracts maximum information (e.g., estimating the mean of a scalar under Gaussian noise with a fixed noise variance).  More generally, $IG(\bm{x})$ is small when the oracle map $\bm{x} \mapsto \bm{s}^\star(\bm{x})$ is nearly constant over the prior support---typically when the likelihood's $\bm{x}$-dependence is weak relative to its $\bm{s}$-dependence.

\subsection{Computability of the ignorance gap}

The ignorance gap is computable from the fixed-design solution alone---no adaptive solver required:
\begin{enumerate}[nosep]
  \item Solve \eqref{eq:fixed} to obtain $\bm{s}^\star_{\mathrm{fixed}}(\bm{c})$.
  \item For a Monte Carlo sample $\{\bm{x}^{(i)}\}_{i=1}^M \sim p_0$:
  \begin{enumerate}[nosep]
    \item Evaluate $\Phi(\bm{x}^{(i)}, \bm{s}^\star_{\mathrm{fixed}}, \bm{c})$.
    \item Solve $V_{\mathrm{oracle}}(\bm{x}^{(i)}; \bm{c}) = \max_{\bm{s}_{1:N}} \Phi(\bm{x}^{(i)}, \bm{s}_{1:N}, \bm{c})$ (an $N d_s$-variable optimization).
    \item Compute $IG(\bm{x}^{(i)})$ as their difference.
  \end{enumerate}
  \item Report $\widehat{\E}[IG]$, $\widehat{\mathrm{std}}[IG]$, $\widehat{\max}[IG]$, $\Pr[IG > \epsilon]$.
\end{enumerate}
If $\widehat{\E}[IG] \ll $ (prior entropy), Corollary~\ref{cor:useless_adapt} says the fixed design is essentially optimal.  This offline check costs $M$ small optimizations and should precede any investment in an adaptive solver.

\paragraph{Diagnostic statistics.}
The distribution of $IG(\bm{x})$ carries more information than its mean:
\begin{center}
\begin{tabular}{@{}ll@{}}
\toprule
Statistic & Interpretation \\
\midrule
$\E[IG]$ & Upper bound on adaptation value (EVPI, in nats) \\
$\mathrm{std}[IG] / \E[IG]$ & Heterogeneity: large $\Rightarrow$ benefit concentrated on certain $\bm{x}$ \\
$\max IG$ & Worst-case suboptimality of fixed design \\
$\Pr(IG > \epsilon)$ & Fraction of state space where adaptation matters \\
\bottomrule
\end{tabular}
\end{center}

%===============================================================================
\section{Adaptive policies}
\label{sec:adaptive}
%===============================================================================

When $\E[IG]$ is not negligible, the fixed design leaves value on the table: by observing intermediate $\bm{y}_k$ and using them to select later $\bm{s}_{k+1}$ we may extract information the fixed schedule misses.  This motivates \emph{adaptive} policies.

\subsection{Definition}

An \textbf{adaptive policy} is a sequence of measurable functions
\begin{equation}
  \pi = \{\pi_1, \pi_2, \ldots, \pi_N\}, \qquad \pi_k: \mathcal{Y}^{k-1} \to \R^{d_s},
\end{equation}
where $\pi_k(\bm{y}_{1:k-1})$ is the sensor action at step $k$, a function of the data observed \emph{before} step $k$.  By convention $\pi_1$ takes no argument.

Equivalently, the history $\bm{y}_{1:k-1}$ enters only through the posterior belief
\begin{equation}
  b_k(\bm{x}) \;=\; p(\bm{x} \mid \bm{y}_{1:k-1}, \bm{s}_{1:k-1}, \bm{c}),
\end{equation}
so we may write $\pi_k: b_k \mapsto \bm{s}_k$---a map from beliefs to sensor parameters.  The belief $b_k$ is the \emph{sufficient statistic} for all future decisions.

\subsection{Contrast with pre-computed sensors}

\begin{center}
\begin{tabular}{@{}lll@{}}
\toprule
Property & Pre-computed (\S\ref{sec:fixed}) & Adaptive (this section) \\
\midrule
$\bm{s}_k$ is & a fixed vector & a function $\pi_k(\bm{y}_{1:k-1}) = \pi_k(b_k)$ \\
Chosen & before any data & after each observation \\
Depends on $\bm{y}_{1:k-1}$? & No & Yes \\
Search space & $\R^{N d_s}$ (finite-dim) & function space (infinite-dim) \\
Benefit & hedges average-case $\bm{x}$ & can target inferred $\hat{\bm{x}}$ \\
Value & $V_{\mathrm{fixed}}$ \eqref{eq:fixed} & $V_{\mathrm{adaptive}} \le V_{\mathrm{fixed}} + \E[IG]$ \\
\bottomrule
\end{tabular}
\end{center}

\begin{remark}
An ``adaptive'' policy is a \emph{function} of the data, not merely a computation performed at run time.  A schedule whose numerical values are computed during data collection but whose \emph{functional form} does not use $\bm{y}_{1:k-1}$ is still a fixed design in this sense.
\end{remark}

%===============================================================================
\section{Finding the optimal adaptive policy: Bellman's equation}
\label{sec:bellman}
%===============================================================================

The adaptive problem
\begin{equation}
  V_{\mathrm{adaptive}}(\bm{c}) \;=\; \max_{\pi} \; \E_{\bm{x}, \bm{y}_{1:N}}\!\bigl[\,\Phi(\bm{x}, \pi(\bm{y}_{1:N-1}), \bm{c})\,\bigr]
  \label{eq:vadapt}
\end{equation}
is a maximization over functions and hence intractable at face value.  The Bellman recursion reduces it to a sequence of single-step maximizations.

\subsection{Value-to-go}

Define the \emph{value-to-go} $W_k(b_k; \bm{c})$: the best expected remaining information gain from step $k$ onward, starting at belief $b_k$.  Boundary values:
\begin{itemize}[nosep]
  \item $W_1(p_0; \bm{c}) = V_{\mathrm{adaptive}}(\bm{c})$---the quantity of interest.
  \item $W_{N+1}(b_{N+1}; \bm{c}) = H(p_0) - H(b_{N+1})$---the accumulated information at the terminal belief.
\end{itemize}
Writing the terminal this way lets $W_{N+1}$ encode the remaining value-to-go when no decisions remain, and $W_1$ equal $\E[\Phi]$ at the prior under the optimal policy.

\subsection{Bellman recursion}

\begin{theorem}[Bellman recursion; see Appendix~\ref{app:bellman} for derivation]
\begin{align}
  W_k(b_k; \bm{c}) &\;=\; \max_{\bm{s}_k}\; \E_{\bm{y}_k \mid b_k, \bm{s}_k, \bm{c}}\!\bigl[\,W_{k+1}(b_{k+1}; \bm{c})\,\bigr], \quad k = N, \ldots, 1, \label{eq:bellman}\\
  W_{N+1}(b_{N+1}; \bm{c}) &\;=\; H(p_0) - H(b_{N+1}), \label{eq:terminal}
\end{align}
where $b_{k+1}$ is the Bayesian posterior update of $b_k$ given $(\bm{s}_k, \bm{y}_k)$.  The optimal policy is
\begin{equation}
  \pi^\star_k(b_k; \bm{c}) \;=\; \argmax_{\bm{s}_k}\; \E_{\bm{y}_k \mid b_k, \bm{s}_k, \bm{c}}\!\bigl[\,W_{k+1}(b_{k+1}; \bm{c})\,\bigr],
  \label{eq:pi_star}
\end{equation}
and $V_{\mathrm{adaptive}}(\bm{c}) = W_1(p_0; \bm{c})$.
\end{theorem}

The equations~\eqref{eq:bellman}--\eqref{eq:terminal} reduce an infinite-dimensional optimization to a sequence of finite-dimensional $\max_{\bm{s}_k}$'s.  The price is that $W_k$ is a function on belief space, which is itself infinite-dimensional in general.  Exact representation is possible in three cases: (i) \emph{finite discrete states}, where $b_k$ lies in a simplex and $W_k$ is piecewise-linear convex (Sondik); (ii) \emph{gridded continuous states with discrete actions and countable observations}, where the reachable belief tree is finite and $W_k$ is tabulated at each reachable $b_k$; (iii) \emph{linear-Gaussian problems}, where $b_k$ is parameterized by its mean and covariance.  Outside these, approximations are required (Section~\ref{sec:relaxations}).

\subsection{Oracle value-to-go and conditional ignorance gap}

Analogous to the unconditional ignorance gap, define
\begin{equation}
  V_k^\star(\bm{x}; \bm{c}) \;=\; \max_{\bm{s}_{k:N}} \Phi(\bm{x}, \bm{s}_{1:k-1}, \bm{s}_{k:N}, \bm{c})
  \label{eq:oracle_vtg}
\end{equation}
(oracle value-to-go at step $k$ given committed $\bm{s}_{1:k-1}$).  Since no policy beats the oracle, $W_k(b_k; \bm{c}) \le \E_{\bm{x}\mid b_k}[V_k^\star(\bm{x}; \bm{c})]$, and the \emph{conditional ignorance gap}
\begin{equation}
  \Gamma_k(b_k; \bm{c}) \;=\; \E_{\bm{x}\mid b_k}[V_k^\star(\bm{x}; \bm{c})] - W_k(b_k; \bm{c}) \;\ge\; 0
\end{equation}
satisfies $\Gamma_1(p_0; \bm{c}) = \E[V_{\mathrm{oracle}}] - V_{\mathrm{adaptive}}$.  Substituting $W_k = \E[V_k^\star] - \Gamma_k$ into \eqref{eq:bellman} yields a DP on $\Gamma_k$ with an explicit exploration--exploitation structure; this \emph{gap DP} is useful for theoretical analysis and is standard in the literature but is not required here.

%===============================================================================
\section{Joint optimization of geometry and adaptive policy}
\label{sec:joint}
%===============================================================================

The full joint problem over the device geometry $\bm{c}$ and the adaptive policy $\pi$ is
\begin{equation}
  \boxed{\;
    \max_{\bm{c}} \; W_1(p_0;\, \bm{c}) \;=\; \max_{\bm{c}} \max_{\pi} \E_{\bm{x}, \bm{y}_{1:N}}\!\bigl[\,\Phi(\bm{x}, \pi(\bm{y}_{1:N-1}), \bm{c})\,\bigr].
  \;}
  \label{eq:joint_bellman}
\end{equation}
The geometry enters every level: the likelihoods $p(\bm{y}_k \mid \bm{x}, \bm{s}_k, \bm{c})$, the posterior-update operator, the terminal $H(b_{N+1})$ via the induced belief trajectory, and hence the oracle schedule $\bm{s}^\star(\bm{x}; \bm{c})$ and the Bellman-optimal policy $\pi^\star(\bm{c})$.

\subsection{Envelope theorem for the outer gradient}
\label{subsec:envelope}

Naive differentiation of \eqref{eq:joint_bellman} would require $\partial\pi^\star/\partial\bm{c}$, a formidable object that is ill-defined at policy-switch boundaries.  The envelope theorem (Danskin) says we don't need it:
\begin{equation}
  \boxed{\;
    \frac{d W_1}{d \bm{c}}\bigg|_{\bm{c} = \bm{c}_0}
    \;=\;
    \frac{\partial}{\partial \bm{c}}\, \E\!\bigl[\,\Phi\!\bigl(\bm{x},\, \pi^\star(\bm{y};\, \bm{c}_0),\, \bm{c}\bigr)\,\bigr]\bigg|_{\bm{c} = \bm{c}_0}.
  \;}
  \label{eq:envelope}
\end{equation}
In words: at the optimal policy for $\bm{c}_0$, the first-order effect of perturbing the policy vanishes (by the first-order optimality condition on $\pi$).  We differentiate $\Phi$ as if the policy were fixed---only the direct effect of $\bm{c}$ on the likelihood and the terminal entropy matter, evaluated at whatever sensor values $\pi^\star(\bm{c}_0)$ happens to emit.

\subsection{Policy iteration}

Equation~\eqref{eq:envelope} enables an alternating policy-iteration scheme:
\begin{enumerate}[nosep, leftmargin=2em]
  \item \textbf{Policy step.}  For current $\bm{c}_t$, solve (or approximate) the Bellman recursion \eqref{eq:bellman} to obtain $\pi^\star(\bm{c}_t)$.
  \item \textbf{Geometry step.}  Take a gradient step on $\bm{c}$ using $\partial\Phi/\partial\bm{c}$ evaluated at $\pi^\star(\bm{c}_t)$ fixed (the RHS of~\eqref{eq:envelope}).
  \item Repeat until convergence.
\end{enumerate}
Step 1 contains the computational burden; step 2 is the same adjoint/sensitivity computation one uses in the fixed-design problem, just with the sensors coming from $\pi^\star$ rather than a constant schedule.

\subsection{Reduction by regime}

The ignorance gap $\E[IG]$ controls how much of this machinery is needed:
\begin{center}
\begin{tabular}{@{}lll@{}}
\toprule
Regime & Joint problem reduces to & Solver \\
\midrule
$\E[IG] \approx 0$ & Fixed design: $\max_{\bm{c}, \bm{s}_{1:N}} \E_{\bm{x}}[\Phi]$ & Single-level gradient descent \\
$\E[IG]$ moderate & Certainty-equivalent policy iteration & Oracle map + filter + adjoint \\
$\E[IG]$ large & Full Bellman DP + policy iteration & Exact DP or parameterized approximation \\
\bottomrule
\end{tabular}
\end{center}
The computability check of Section~\ref{sec:ig} decides which row applies before any solver is implemented.

\subsection{The added value of joint optimization}
\label{subsec:value_of_joint}

The decoupled program
\begin{center}
``\emph{Pick $\bm{c}$ to maximize $\E[IG](\bm{c})$ (since large EVPI promises adaptive benefit).  Separately, at run time, execute an adaptive policy.}''
\end{center}
is \emph{not} equivalent to the joint optimization \eqref{eq:joint_bellman}.  Three facts together make joint optimization a distinct and necessary step:

\begin{enumerate}[nosep, leftmargin=2em]
\item $\E[IG](\bm{c})$ and $V_{\mathrm{adaptive}}(\bm{c}) - V_{\mathrm{fixed}}(\bm{c})$ are \emph{different functions of $\bm{c}$}.  A $\bm{c}$ maximizing one need not maximize the other.
\item The ratio $\bigl(V_{\mathrm{adaptive}} - V_{\mathrm{fixed}}\bigr) / \E[IG]$ (the \emph{EVPI saturation fraction}) varies substantially with $\bm{c}$ and is typically well below $1$.
\item The $\bm{c}$ maximizing $V_{\mathrm{adaptive}}$ generally differs from the $\bm{c}$ maximizing $V_{\mathrm{fixed}}$.
\end{enumerate}

\paragraph{Radar case study, within-Gaussian-family comparison (see Section~\ref{sec:examples}).}
For the $K = 16$, $N = 4$ benchmark at $p_{\mathrm{D}} = 0.9$, $p_{\mathrm{FA}} = 0.05$, restricting attention to the Gaussian-beam family $G_w(\alpha) = \exp(-\alpha^2/(2w^2))$ and comparing the two extremes against the in-family optimum $w^\star = 0.60$~rad:
\begin{center}
\begin{tabular}{@{}lcccc@{}}
\toprule
geometry $\bm{c}$ & $\E[IG]$ (EVPI) & $V_{\mathrm{adaptive}}$ & realized & saturation \\
\midrule
Gaussian $w = 0.05$ (small-$w$ extreme)   & $2.08$ nats & $0.62$ nats           & $0.12$ nats          & $5.8\%$           \\
Gaussian $w = \pi$\phantom{.05} (large-$w$ extreme) & $0.09$ nats & $0.12$ nats           & $0.002$ nats         & $2.4\%$           \\
Gaussian $w = 0.60$ (in-family $w^\star$) & $0.64$ nats & $\mathbf{0.87}$ nats  & $0.11$ nats          & $\mathbf{17.0\%}$ \\
\bottomrule
\end{tabular}
\end{center}
The small-$w$ extreme has the \emph{largest} $\E[IG]$ (by a factor of $\sim\!3$ over $w^\star$), yet only the second-largest $V_{\mathrm{adaptive}}$ and a saturation of just $5.8\%$.  The large-$w$ extreme has both the smallest $\E[IG]$ and the smallest $V_{\mathrm{adaptive}}$.  The in-family optimum $w = 0.60$ \emph{dominates} both extremes on $V_{\mathrm{adaptive}}$ despite having neither the largest $\E[IG]$ nor the largest realized gain.  Routing the geometry choice through EVPI alone would have picked the small-$w$ design; the joint DP picks $w = 0.60$.

\paragraph{Practical prescription.}
\begin{enumerate}[nosep, leftmargin=2em]
  \item Compute $\E[IG]$ at a few candidate $\bm{c}$.  If uniformly $\ll$ prior entropy, \eqref{eq:fixed} suffices (Corollary~\ref{cor:useless_adapt}).
  \item Otherwise run the joint $(\bm{c}, \pi)$ optimization \eqref{eq:joint_bellman} via policy iteration with envelope-theorem geometry gradients.  The EVPI diagnostic determines \emph{whether} to engage this loop; the joint DP determines \emph{where} its optimum lies.
\end{enumerate}
EVPI and joint-DP serve complementary but non-interchangeable roles.

%===============================================================================
\section{A tractable case study: adaptive radar beam search}
\label{sec:examples}
%===============================================================================

We conclude with a concrete problem that exercises every ingredient of the framework: $\E[IG]$ is demonstrably large, the geometry $\bm{c}$ exerts first-order control on achievable performance, and the Bellman recursion \eqref{eq:bellman} is exactly solvable.  The policy-iteration scheme of Section~\ref{sec:joint} becomes a concrete algorithm.  The problem is deliberately small and canonical so that every claim can be verified numerically.

\subsection{Problem setup}
\label{subsec:radar_setup}

\paragraph{State.}  A single target occupies one of $K$ discrete cells at distinct angular positions.  We use $K = 16$ cells equally spaced on a circular arc at fixed range from the antenna, so cell $z \in \{1, \ldots, K\}$ has bearing $\varphi(z) = 2\pi(z-1)/K$.  The prior on $z$ is uniform.

\paragraph{Actions.}  At each step $k = 1, \ldots, N$ (we use $N = 4$), the operator steers the antenna to a pointing direction $\theta_k \in \Theta = \{\theta_1, \ldots, \theta_M\}$ taken from a finite set of candidate pointings (we take $M = K$).

\paragraph{Observations.}  Each pointing returns a binary detection $y_k \in \{0, 1\}$ with Bernoulli likelihood
\begin{equation}
  \Pr(y_k = 1 \mid z, \theta_k; \bm{c}) \;=\; p_{\mathrm{FA}} + (p_{\mathrm{D}} - p_{\mathrm{FA}})\, G_{\bm{c}}\!\bigl(\theta_k - \varphi(z)\bigr),
  \label{eq:radar_likelihood}
\end{equation}
where $G_{\bm{c}}(\alpha) \in [0, 1]$ is the normalized antenna gain pattern at angular offset $\alpha$ from the main-lobe direction.

\paragraph{Geometry.}  The design vector $\bm{c}$ is the full set of antenna-array parameters: a phased array of $N_e$ elements at positions $\{\bm{r}_n\}$ driven by complex excitations $\{a_n\}$.  The far-field array factor is
\begin{equation}
  F_{\bm{c}}(\alpha) \;=\; \sum_{n=1}^{N_e} a_n \,\exp\!\bigl(\mathrm{i}\, k_0\, \bm{r}_n \cdot \hat{\bm{\alpha}}\bigr), \qquad k_0 = 2\pi/\lambda,
  \label{eq:array_factor}
\end{equation}
and $G_{\bm{c}}(\alpha) = |F_{\bm{c}}(\alpha)|^2 / \max_{\alpha'}|F_{\bm{c}}(\alpha')|^2$.  The design vector is $\bm{c} = (\{\bm{r}_n\}, \{a_n\})$ with real dimension $d_c \sim 2$--$4 N_e$ depending on what is free.  Typical $N_e \in \{8, \ldots, 64\}$ gives $d_c$ in the tens to low hundreds.

\paragraph{Reward.}  For this discrete problem the reward is
\begin{equation}
  \Phi(z, \theta_{1:N}, \bm{c}) \;=\; \ln K - H(z \mid y_{1:N}, \theta_{1:N}, \bm{c})
\end{equation}
---the discrete specialization of Definition~\ref{def:phi} with $H(p_0) = \ln K$ (uniform prior over $K$ cells).  Thus $\Phi \in [0, \ln K]$, with $\Phi = \ln K$ at full resolution.

\subsection{The ignorance gap is large}

\paragraph{Oracle.}  Knowing $z$, the oracle always points $\theta_k^\star = \varphi(z)$ (or the nearest element of $\Theta$), receiving an on-target Bernoulli with probability $p_{\mathrm{D}}$ every step.  For $p_{\mathrm{D}} \gg p_{\mathrm{FA}}$ and $N \ge \log_2 K$ (saturated at $K = 16$, $N = 4$), the posterior is essentially fully resolved: $V_{\mathrm{oracle}}(z) \approx \ln K$.

\paragraph{Fixed design (narrow-beam family).}  Assume geometry $\bm{c}$ produces a single-main-lobe, directional gain with main-lobe width $\Delta\theta_{\mathrm{ML}}(\bm{c}) \lesssim 2\pi/K$ (one cell per main lobe).  With $N < K$ narrow beams, at most $N$ distinct cells are probed.  Conditioning on whether the target lies inside one of those $N$ cells (probability $N/K$) or outside (probability $(K-N)/K$):
\begin{equation}
  V_{\mathrm{fixed}} \;\approx\; \ln K - \frac{K - N}{K}\ln(K - N).
  \label{eq:radar_vfixed}
\end{equation}

\paragraph{EVPI.}  Combining,
\begin{equation}
  \boxed{\;\E[IG] \;=\; V_{\mathrm{oracle}} - V_{\mathrm{fixed}} \;\approx\; \frac{K - N}{K}\ln(K - N) \;\text{ nats},\;}
  \label{eq:radar_ig}
\end{equation}
which scales as $\sim \ln K$ for $K \gg N$ and vanishes as $N \to K$.  At $K = 16$, $N = 4$, $\E[IG] \approx (12/16)\ln 12 \approx 1.86$ nats---a substantial EVPI upper bound whose size, however, does \emph{not} imply a large adaptive advantage (cf.\ Remark~\ref{rem:evpi_necessary_not_sufficient}).

\subsection{Geometry has first-order impact}

The design $\bm{c}$ shapes every feature of $G_{\bm{c}}$ independently---main-lobe width, sidelobe structure, nulls, grating lobes---and each feeds into the Bernoulli likelihood~\eqref{eq:radar_likelihood} in a distinct way.  Main-lobe width (aperture size), first sidelobe level (amplitude taper), grating-lobe behavior (element spacing), and null placement (complex excitations) all have first-order impact on both $V_{\mathrm{adaptive}}$ and $V_{\mathrm{fixed}}$, and they interact.  The joint optimum over $(\bm{c}, \pi)$ solves a $d_c$-dimensional problem whose objective is modulated by the adaptive policy that will run on top.

\paragraph{Rough sizing.}
Binary observations limit single-step information to $\ln 2$ nats, so full $K$-cell resolution requires $N \ge \log_2 K$.  For a flat ``in-beam / out-of-beam'' idealization with $m$ cells per main lobe, single-measurement mutual information is maximized at $m^\star \approx K/2$, giving main-lobe width
\begin{equation}
  \Delta\theta_{\mathrm{ML}}^\star \;\sim\; \pi \qquad \text{(half the ring).}
  \label{eq:radar_width_scaling}
\end{equation}
A key fact: $N$ contiguous-arc beams produce at most $2N$ arc boundaries and hence $2N$ distinguishable segments of $K/(2N)$ cells each.  Adaptive policies can bisect the current posterior each step---a binary search reaching single-cell resolution in $\log_2 K$ steps---while fixed designs cannot.

Two regimes at opposite extremes of the beam-width knob:

\begin{itemize}[nosep, leftmargin=*]
\item \emph{Narrow} ($m \approx 1$, $\Delta\theta_{\mathrm{ML}} \lesssim 2\pi/K$).  Each measurement probes one cell; fixed and adaptive achieve $V_{\mathrm{fixed}} = V_{\mathrm{adaptive}} \approx \ln K - \frac{K-N}{K}\ln(K-N) \approx 0.91$ nats at $K=16, N=4$.  EVPI $= 1.86$ nats; adaptive advantage $= 0$.  \textbf{Large EVPI, zero realized adaptive advantage}---the canonical counter-example to the ``large EVPI implies valuable adaptation'' fallacy.

\item \emph{Wide} ($m \approx K/2$, $\Delta\theta_{\mathrm{ML}} \sim \pi$).  Fixed pointings give at most $\ln(2N) \approx 2.08$ nats; adaptive binary search gives $\ln K \approx 2.77$ nats.  Adaptive advantage $= \ln(K/(2N)) = \ln 2 \approx 0.69$ nats \emph{equals} $\E[IG]$ here---EVPI is \textbf{tight} and adaptation realizes the full gap.
\end{itemize}

The narrow regime maximizes the \emph{computable} EVPI upper bound but fails to saturate it; the wide regime has smaller $\E[IG]$ but saturates it fully.  Joint optimization favors the wide design: both $V_{\mathrm{fixed}}^\star = \ln(2N)$ and $V_{\mathrm{adaptive}}^\star = \ln K$ are maximized by $m \approx K/2$.

\subsection{Full Bellman DP is tractable}

This is a finite-state, finite-action, finite-observation POMDP with short horizon.  For the terminal reward $\Phi = \ln K - H(b_{N+1})$ (nonlinear in $b$), the Sondik PWLC representation does not apply directly, but direct backward induction on the reachable belief tree does: from the uniform prior, binary observations under each of $K$ actions produce a tree of at most $(2K)^N$ leaves.  At $K = 16, N = 4$ this is $\sim 10^6$ leaves---trivially tabulable.  Backward induction with memoization runs in seconds on a laptop.  Symmetry reduction by the dihedral $D_K$ group ($|G| = 32$) further reduces the work by up to an order of magnitude.

\subsection{Policy iteration for the joint optimum}

With exact inner DP available, the envelope-theorem scheme of Section~\ref{sec:joint} becomes a concrete algorithm:

\begin{enumerate}[nosep, leftmargin=2em]
  \item \textbf{Policy step.}  For current $\bm{c}_t$, solve the finite POMDP exactly to obtain $\pi^\star(\bm{c}_t)$ as a depth-$N$ policy tree: every branch indexed by $y_{1:k-1}$ maps to an action $\theta_k$.
  \item \textbf{Geometry step.}  $\Phi$ depends on $\bm{c}$ analytically through \eqref{eq:radar_likelihood} and \eqref{eq:array_factor}.  Compute
  \begin{equation}
    \nabla_{\bm{c}} W_1\big|_{\bm{c}_t} \;=\; \E_{z, y_{1:N} \sim \pi^\star(\bm{c}_t)}\!\bigl[\nabla_{\bm{c}} \Phi\!\bigl(z, \pi^\star(y_{1:N}; \bm{c}_t), \bm{c}\bigr)\bigr]_{\bm{c} = \bm{c}_t}
    \label{eq:radar_env}
  \end{equation}
  by Monte Carlo over $(z, y_{1:N})$ trajectories under the frozen $\pi^\star(\bm{c}_t)$.  The envelope theorem means no policy Jacobian $d\pi^\star / d\bm{c}$ is needed.
  \item \textbf{Gradient step.}  Update $\bm{c}_{t+1} = \mathrm{Proj}_{\mathcal{C}}(\bm{c}_t + \eta \nabla_{\bm{c}} W_1)$ and repeat.
\end{enumerate}

Because step 1 is \emph{exact}, the gradient direction \eqref{eq:radar_env} is exact---only Monte Carlo variance over trajectories shrinks as $|\text{samples}|^{-1/2}$.  Convergence is governed by the smoothness of $W_1(\cdot)$ in $\bm{c}$: continuous except on codimension-one policy-switch boundaries, through which a subgradient substitutes for the gradient.

\subsection{Quantitative demonstration of the necessary-but-not-sufficient property}

In the companion numerics document (with code in the same directory), we sweep $\bm{c}$ over the Gaussian beamwidth $w \in (0, \pi]$ and compute exact $V_{\mathrm{oracle}}$, $V_{\mathrm{fixed}}$, $V_{\mathrm{adaptive}}$, $\E[IG]$ at each $w$.  The in-family optimum for $V_{\mathrm{adaptive}}$ is $w^\star = 0.60$~rad with $V_{\mathrm{adaptive}}^\star = 0.87$ nats; the EVPI-maximizing $w$ is the small-$w$ extreme with $V_{\mathrm{adaptive}} = 0.62$ nats---a $1.4\times$ ratio.  Within the top-hat family (integer $m \in \{1, \ldots, K\}$), the in-family optimum is $m^\star = 9$ with $V_{\mathrm{adaptive}}^\star = 1.70$ nats, while EVPI-maximizing $m$ gives $V_{\mathrm{adaptive}} = 0.62$ nats---a $2.8\times$ ratio.  The largest-EVPI designs in each family are the \emph{smallest}-$V_{\mathrm{adaptive}}$ designs, quantitatively confirming Remark~\ref{rem:evpi_necessary_not_sufficient} (and the analog of the joint-optimum finding in \S\ref{subsec:value_of_joint}).

\subsection{Adjacent tractable testbeds}

The same template covers several other finite-POMDP problems with meaningful geometry:
\begin{itemize}[nosep]
  \item \emph{Sequential material identification by spectroscopy}: finite hypothesis set, excitation wavelengths, quantized absorbance; $\bm{c}$ = spectrometer linewidth and grid.
  \item \emph{Adaptive compressed sensing with known sparsity}: discrete support set, dictionary rows as actions, thresholded projections as observations; $\bm{c}$ = dictionary design.
  \item \emph{Sequential medical-imaging slicing}: discrete lesion-location hypotheses, scan parameters as actions, quantized intensities as observations; $\bm{c}$ = field-of-view and gradient design.
\end{itemize}

%===============================================================================
\section{When the joint DP is intractable: a hierarchy of principled relaxations}
\label{sec:relaxations}
%===============================================================================

The radar case study of Section~\ref{sec:examples} is engineered for exact Bellman tractability: finite state, finite observation alphabet, finite action grid, short horizon.  Many sensor-design problems of practical interest violate at least one of these.  When exact Bellman is beyond reach, the framework of Sections~\ref{sec:bellman}--\ref{sec:joint} remains the correct organizing picture; only the inner optimization needs approximation.  This section catalogs principled relaxations, ordered by structure preserved.

\subsection{Sources of intractability}

Three structural barriers recur:
\begin{enumerate}[leftmargin=2em, nosep]
  \item \textbf{Belief-space size.}  Continuous $\bm{x}$ yields an infinite-dimensional belief; one must discretize (grid), project onto a parametric family (e.g., Gaussian), or carry a particle representation.
  \item \textbf{Action-space size.}  Continuous $\bm{s}$ (beamwidth, amplitude taper, phase profile) turns the inner $\max_{\bm{s}_k}$ into a nonconvex continuous optimization.
  \item \textbf{Per-step cost.}  Physics-based forward models (FDFD, FDTD, RCWA) make each belief-tree node expensive.
\end{enumerate}

Any one collapses the exact DP; in realistic photonic / metasurface settings all three bite.

\subsection{Relaxation path 1: neural adaptive policies (deep RL)}
\label{subsec:deep_rl}

Parametrize the policy as a neural network $\pi_\theta$ and train $\theta$ jointly with $\bm{c}$ by policy-gradient / REINFORCE / pathwise-gradient methods.  The joint objective is
\begin{equation}
  \max_{\bm{c},\,\theta}\; \E_{\bm{x}, \bm{y}_{1:N}}\bigl[\,\Phi(\bm{x},\bm{s}_{1:N}^{\pi_\theta},\bm{c})\,\bigr].
  \label{eq:deep_rl_joint}
\end{equation}

\paragraph{What this preserves.}  Multi-step reasoning via NN memory of the history; arbitrary action geometry; pathwise $\partial\Phi/\partial\bm{c}$ if the forward model is differentiable.  Modern end-to-end optical learning (Sitzmann et al., 2018; Chang \& Wetzstein, 2019) fits this template with $\pi_\theta$ collapsed to a one-step reconstruction network.

\paragraph{What this gives up.}  Optimality certificates: $\pi_\theta$ is an approximator, not a Bellman solver.  The envelope theorem \eqref{eq:envelope} applies cleanly only at an interior optimum of $\theta$, rarely attained in practice.

\paragraph{Failure mode: hardware-light fixed points.}  A familiar pathology of end-to-end optical learning: $\pi_\theta$ can absorb the training signal and compensate for a poor $\bm{c}$ by investing in policy complexity, leaving $\bm{c}$ near initialization.  The mechanism is gradient asymmetry---$\partial\Phi/\partial\theta$ propagates through cheap high-rank NN backprop, $\partial\Phi/\partial\bm{c}$ through an expensive low-rank forward-model adjoint.  In regimes where geometry has first-order impact on $V_{\mathrm{adaptive}}$ (Section~\ref{sec:examples} quantifies this as $\sim\!2.8\times$ in the radar case), this failure mode silently returns a strictly worse answer than the bi-level paths below.  Standard countermeasures: (i) capacity control on $\pi_\theta$; (ii) two-timescale updates with $\bm{c}$ faster than $\theta$; (iii) warm-start $\theta$ from a prescribed information-theoretic rule (2b)/(2c) below; (iv) regularizers penalizing ``easy'' $\bm{c}$.  In our experience (iii) is most reliable.

\subsection{Relaxation path 2: bi-level optimization with a prescribed inner policy}
\label{subsec:bilevel}

A cheaper, more interpretable route prescribes $\pi$ in closed form as an information-theoretic rule and keeps the outer envelope-theorem loop of Section~\ref{sec:joint}:
\begin{equation}
  \max_{\bm{c}}\; \E\bigl[\,\Phi(\bm{x},\,\bm{s}_{1:N}[b_{\cdot};\bm{c}],\,\bm{c})\,\bigr]
  \qquad\text{subject to } \bm{s}_k = \mathsf{R}_k(b_k;\bm{c}).
  \label{eq:bilevel_general}
\end{equation}
The outer geometry gradient flows through the implicit $\bm{c}$-dependence of the inner $\argmax$ via the implicit function theorem (IFT)---exact even when the inner rule is not Bellman-optimal.  The rule $\mathsf{R}_k$ determines how much Bellman structure survives.

\subsubsection*{(2a) Oracle certainty-equivalent control}

At step $k$, form $\hat{\bm{x}}_k = \argmax_{\bm{x}} b_k$ (MAP) or $\hat{\bm{x}}_k = \E[\bm{x}\mid b_k]$, then play the first step of the oracle schedule at $\hat{\bm{x}}_k$:
\begin{equation}
  \bm{s}_k^{\mathrm{CE}} = \bigl(\bm{s}^\star(\hat{\bm{x}}_k;\bm{c})\bigr)_1,
  \qquad
  \bm{s}^\star(\bm{x};\bm{c}) = \argmax_{\bm{s}_{1:N-k+1}}\,\Phi(\bm{x},\,\bm{s}_{1:N-k+1},\,\bm{c}).
  \label{eq:ce_rule}
\end{equation}
Receding-horizon, non-myopic, but assumes $\hat{\bm{x}}_k$ is correct.  Characteristic failure: narrow-beam geometries where $\bm{s}^\star(\bm{x})$ collapses to a single pointing at $\bm{x}$ induce exploration deficits (point-and-verify with a wrong guess).  Computationally, one oracle solve per step---if $V_{\mathrm{oracle}}$ is the bottleneck, this relaxation does not escape it.

\subsubsection*{(2b) Myopic one-step information-maximization}

Choose each $\bm{s}_k$ to maximize the expected posterior entropy reduction from the next observation alone:
\begin{equation}
  \bm{s}_k^{\mathrm{myopic}} \;=\; \argmax_{\bm{s}_k}\; I(\bm{x};\, y_k \mid b_k, \bm{s}_k, \bm{c}).
  \label{eq:myopic_rule}
\end{equation}
No look-ahead beyond one step.  Greedy information gain is $(1-1/e)$-competitive against the clairvoyant DP optimum for submodular objectives (Krause \& Guestrin, 2008; Golovin \& Krause, 2011), the default baseline in sequential design.

\subsubsection*{(2c) Gaussian-posterior / Fisher-trace surrogate with recursive inference}

In the Gaussian-posterior / BvM regime of Section~\ref{subsec:fisher_specialization}, $\Phi(\bm{x},\bm{s},\bm{c}) \approx \tfrac{1}{2}\ln\det \bm{J}(\bm{x},\bm{s},\bm{c}) + \mathrm{const}$, where $\bm{J}$ is the accumulated Fisher information.  A stack of standard approximations on top of (2b) yields the cheapest member of the family:
\begin{equation}
  \bm{s}_k^{\mathrm{BFIM}} \;=\; \argmax_{\bm{s}_k}\;\tr \bm{J}\bigl(\bar{\bm{x}}_k,\,\bm{s}_k,\,\bm{c}\bigr),
  \qquad
  \bar{\bm{x}}_k = \E[\bm{x}\mid b_k],
  \label{eq:bfim_rule}
\end{equation}
coupled with an extended Kalman filter for the belief update.  The three stacked simplifications relative to (2b) are:
\begin{enumerate}[label=(\roman*), leftmargin=2em, nosep]
  \item $\tr$ in place of $\ln\det$ (A-optimal, cheaper; loses scale-invariance).
  \item Evaluate $\bm{J}$ at $\bar{\bm{x}}_k$ instead of $\E_{\bm{x}\mid b_k}[\bm{J}(\bm{x},\cdot)]$ (Laplace-type approximation).
  \item EKF for $b_k$ (Gaussian posterior, linearized updates).
\end{enumerate}
Each individually is standard.  Stacked, they reduce the inner step to an analytic $\max$-of-trace at a single $\bm{x}$-point, with a closed-form IFT adjoint for $\bm{c}$-gradients.

\subsection{The hierarchy}

\begin{table}[h]
\centering
\caption{Tractability hierarchy for adaptive sensor design.  Every row is a bi-level problem in $\bm{c}$; only the fidelity of the inner optimization changes.  The exact DP is the gold standard the radar case study of Section~\ref{sec:examples} exhibits.}
\label{tab:hierarchy}
\small
\begin{tabular}{@{}lp{3.4cm}p{2.5cm}p{3.1cm}p{2.5cm}@{}}
\toprule
& inner step per $\bm{s}_k$ & look-ahead & posterior handling & guarantee \\
\midrule
Exact Bellman DP & full belief backup & $N$-step & exact & joint-optimal \\
Deep RL (\S\ref{subsec:deep_rl}) & one NN forward pass & $\le N$, implicit & NN-internal & approximate, black-box \\
Oracle CE (2a) & oracle at $\hat{\bm{x}}_k$ & $(N{-}k{+}1)$-step CE & point estimate & none; under-explores \\
Myopic info-gain (2b) & $\argmax I(\bm{x}; y_k\mid b_k)$ & 1-step & full posterior & $(1-1/e)$ via submodularity \\
Fisher-trace (2c) & $\argmax \tr\bm{J}(\bar{\bm{x}}_k)$ & 1-step & Gaussian / EKF & asymptotically exact under BvM \\
\bottomrule
\end{tabular}
\end{table}

Two observations.  First, \emph{every row} is a bi-level problem in $\bm{c}$: the inner step is a function of $(b_k, \bm{c})$, and the geometry gradient flows through by IFT.  The joint $(\bm{c}, \pi)$ framework of Section~\ref{sec:joint} is structurally identical across rows; only inner fidelity changes.  Second, path 1 (deep RL) and paths 2a--2c are not mutually exclusive: a deep-RL policy can be warm-started from a prescribed rule, combining the sample efficiency of (2b)--(2c) with the flexibility of (1).

\subsection{Transition: the photonic sensor problem as an instance of (2c)}
\label{subsec:transition_photonic}

A companion body of work develops a worked instance of path (2c) applied to a 2D photonic sensor: a cross-waveguide device whose permittivity distribution is the geometry variable $\bm{c}$, whose sensor parameters are input-port phases, and whose observations are port powers.  Each of the three bulleted simplifications (i)--(iii) above is justified by the operating regime:
\begin{itemize}[leftmargin=2em, nosep]
  \item The state $\bm{x}$ is a small refractive-index perturbation ($\Delta n \in [10^{-5}, 10^{-4}]$), so the posterior remains unimodal and near-Gaussian---justifying (iii).
  \item The Fisher information $\bm{J}(\bm{x},\bm{s},\bm{c})$ varies slowly in $\bm{x}$ over the posterior support---justifying (ii).
  \item $\tr\bm{J}$ admits an analytically differentiable inner $\argmax$---justifying (i).
\end{itemize}

The photonic problem is not an alternative \emph{framework}; it is an instantiation of the joint-$(\bm{c}, \pi)$ framework of this paper at row (2c), with each approximation justified by the physics of its operating regime.  The radar case study of Section~\ref{sec:examples} and the photonic case study are two points on the same methodological spectrum: exact DP where tractable, (2c) where not.  The message of both is the same: \emph{joint} optimization of $(\bm{c}, \pi)$ is the correct object; the inner solver depends on what is affordable.

%===============================================================================
\section*{References}
%===============================================================================

\begin{enumerate}[label={[\arabic*]}]
  \item V.~V.~Fedorov, \emph{Theory of Optimal Experiments}.  Academic Press, 1972.
  \item K.~Chaloner and I.~Verdinelli, ``Bayesian experimental design: A review,'' \emph{Statistical Science}, vol.~10, no.~3, pp.~273--304, 1995.
  \item D.~V.~Lindley, ``On a measure of the information provided by an experiment,'' \emph{Ann.\ Math.\ Statist.}, vol.~27, no.~4, pp.~986--1005, 1956.
  \item R.~Bellman, \emph{Dynamic Programming}.  Princeton University Press, 1957.
  \item J.-M.~Bernardo, ``Expected information as expected utility,'' \emph{Ann.\ Statist.}, vol.~7, no.~3, pp.~686--690, 1979.
  \item R.~A.~Howard, ``Information value theory,'' \emph{IEEE Trans.\ Systems Science and Cybernetics}, vol.~2, no.~1, pp.~22--26, 1966.
  \item A.~Krause and C.~Guestrin, ``Near-optimal observation selection using submodular functions,'' \emph{Proc.\ AAAI}, 2007.
  \item D.~Golovin and A.~Krause, ``Adaptive submodularity: Theory and applications in active learning and stochastic optimization,'' \emph{J.\ Artif.\ Intell.\ Res.}, vol.~42, pp.~427--486, 2011.
  \item J.~Haupt, R.~Castro, and R.~Nowak, ``Distilled sensing: Adaptive sampling for sparse detection and estimation,'' \emph{IEEE Trans.\ Information Theory}, vol.~57, no.~9, pp.~6222--6235, 2011.
  \item M.~L.~Hernandez, ``Optimal sensor management for target tracking applications,'' Ph.D.\ thesis, University of Cambridge, 2004.
  \item V.~Sitzmann et al., ``End-to-end optimization of optics and image processing for achromatic extended depth of field and super-resolution imaging,'' \emph{ACM Trans.\ Graph.}, vol.~37, no.~4, 2018.
  \item J.~Chang and G.~Wetzstein, ``Deep optics for monocular depth estimation and 3D object detection,'' \emph{Proc.\ ICCV}, 2019.
  \item S.~P.~Boyd and L.~Vandenberghe, \emph{Convex Optimization}.  Cambridge University Press, 2004.
\end{enumerate}

\appendix

%===============================================================================
\section{Identities for the pointwise information gain $\Phi$}
\label{app:phi_identities}
%===============================================================================

Three equivalent forms of $\Phi(\bm{x}, \bm{s}, \bm{c})$ in Definition~\ref{def:phi}:

\paragraph{Entropy-reduction form (primary).}
\begin{equation}
  \Phi \;=\; H(p_0) - \E_{\bm{y} \sim p(\cdot\mid\bm{x},\bm{s},\bm{c})}[H(p(\cdot\mid\bm{y},\bm{s},\bm{c}))].
\end{equation}

\paragraph{Expected-KL form.}
\begin{equation}
  \Phi \;=\; \E_{\bm{y} \sim p(\cdot\mid\bm{x},\bm{s},\bm{c})}\!\bigl[\KL\bigl(p(\cdot\mid\bm{y},\bm{s},\bm{c})\,\big\|\, p_0\bigr)\bigr].
\end{equation}

\paragraph{Log-likelihood-ratio form.}
\begin{equation}
  \Phi \;=\; \E_{\bm{y} \sim p(\cdot\mid\bm{x},\bm{s},\bm{c})}\!\left[ \int p(\bm{x}'\mid\bm{y},\bm{s},\bm{c}) \ln \frac{p(\bm{x}'\mid\bm{y},\bm{s},\bm{c})}{p_0(\bm{x}')} d\bm{x}' \right].
\end{equation}

Equivalence follows from $\KL(p\,\|\,p_0) = \int p \ln(p/p_0) = -H(p) - \int p\ln p_0$ combined with $\E_{\bm{y}\mid\bm{x}} \int p(\bm{x}'\mid\bm{y}) \ln p_0(\bm{x}') d\bm{x}' = \E_{\bm{y}\mid\bm{x}} \E_{\bm{x}'\mid\bm{y}}[\ln p_0(\bm{x}')]$; expanding and simplifying gives the stated identities.

The prior-average identity $\E_{\bm{x}}[\Phi] = I(\bm{x};\bm{y})$ (Theorem~\ref{thm:phi_is_mi}) is fundamental to the framework.  It means: what the oracle / ignorance-gap / joint-DP machinery does pointwise, the standard mutual-information framework does in aggregate---with the pointwise version strictly richer because it supports oracle conditioning on $\bm{x}$.

%===============================================================================
\section{Derivation of Bellman's equation}
\label{app:bellman}
%===============================================================================

We derive equations~\eqref{eq:bellman}--\eqref{eq:terminal} from first principles.  The argument is general to any finite-horizon sequential decision problem.

\subsection{Problem setup}

Recall the ingredients:
\begin{itemize}[nosep]
  \item An unknown state $\bm{x} \sim p_0$.
  \item Horizon $N$ with steps $k = 1, 2, \ldots, N$.
  \item At step $k$, an action $\bm{s}_k$ is chosen; then $\bm{y}_k \sim p(\cdot \mid \bm{x}, \bm{s}_k, \bm{c})$ is observed.
  \item An adaptive policy $\pi = (\pi_1, \ldots, \pi_N)$ with $\pi_k: \bm{y}_{1:k-1} \mapsto \bm{s}_k$.
  \item A terminal reward $R(b_{N+1}) = H(p_0) - H(b_{N+1})$ accrued at the end.
\end{itemize}

Under policy $\pi$ and initial prior $p_0$, the expected reward is
\begin{equation}
  V^{\pi}(\bm{c}) \;=\; \E_{\bm{x}, \bm{y}_{1:N}}\!\bigl[R(b_{N+1})\bigr].
\end{equation}
By construction, $V^{\pi}(\bm{c}) = \E_{\bm{x}, \bm{y}_{1:N}}[\Phi(\bm{x}, \pi(\bm{y}_{1:N-1}), \bm{c})]$: terminal-entropy value-to-go equals the expected pointwise information gain.

The goal:
\begin{equation}
  V_{\mathrm{adaptive}}(\bm{c}) \;=\; \max_\pi V^{\pi}(\bm{c}).
\end{equation}

\subsection{Sufficient statistic: the belief}

All information in $\bm{y}_{1:k-1}$ relevant to future decisions is carried by the posterior
\begin{equation}
  b_k(\bm{x}) \;=\; p(\bm{x} \mid \bm{y}_{1:k-1}, \bm{s}_{1:k-1}, \bm{c}).
\end{equation}
Two histories yielding the same $b_k$ are indistinguishable going forward, so without loss we search over policies $\pi_k: b_k \mapsto \bm{s}_k$.

Bayesian updating: given $(b_k, \bm{s}_k, \bm{y}_k)$,
\begin{equation}
  b_{k+1}(\bm{x}) \;=\; \frac{p(\bm{y}_k \mid \bm{x}, \bm{s}_k, \bm{c})\, b_k(\bm{x})}{\int p(\bm{y}_k \mid \bm{x}', \bm{s}_k, \bm{c})\, b_k(\bm{x}')\, d\bm{x}'}.
\end{equation}

\subsection{Value-to-go}

Define
\begin{equation}
  W_k(b_k; \bm{c}) \;=\; \max_{\pi_k, \ldots, \pi_N}\; \E\!\bigl[R(b_{N+1}) \mid b_k, \pi_{k:N}\bigr].
  \label{eq:app:vtg}
\end{equation}
Boundary cases:
\begin{itemize}[nosep]
  \item $k = 1$: $W_1(p_0; \bm{c}) = V_{\mathrm{adaptive}}(\bm{c})$.
  \item $k = N + 1$: no decisions remain; $W_{N+1}(b_{N+1}; \bm{c}) = R(b_{N+1}) = H(p_0) - H(b_{N+1})$.
\end{itemize}

\subsection{Peeling off one step}

Split the outer max:
\begin{equation}
  W_k(b_k; \bm{c}) \;=\; \max_{\bm{s}_k}\; \max_{\pi_{k+1}, \ldots, \pi_N}\; \E[R \mid b_k, \bm{s}_k, \pi_{k+1:N}].
\end{equation}
Apply the tower property, conditioning on $\bm{y}_k$ (which determines $b_{k+1}$ given $(b_k, \bm{s}_k)$):
\begin{equation}
  \E[R \mid b_k, \bm{s}_k, \pi_{k+1:N}] \;=\; \E_{\bm{y}_k \mid b_k, \bm{s}_k}\!\bigl[\,\E[R \mid b_{k+1}, \pi_{k+1:N}]\,\bigr].
\end{equation}
Pushing $\max_{\pi_{k+1:N}}$ inside $\E_{\bm{y}_k}$ (justified in the next subsection):
\begin{equation}
  W_k(b_k; \bm{c}) \;=\; \max_{\bm{s}_k}\; \E_{\bm{y}_k \mid b_k, \bm{s}_k}\!\bigl[\max_{\pi_{k+1:N}} \E[R \mid b_{k+1}, \pi_{k+1:N}]\bigr] \;=\; \max_{\bm{s}_k}\; \E_{\bm{y}_k \mid b_k, \bm{s}_k}\!\bigl[W_{k+1}(b_{k+1}; \bm{c})\bigr].
  \label{eq:app:bellman}
\end{equation}
The boxed equation~\eqref{eq:bellman} is \eqref{eq:app:bellman}.

\subsection{Why the push-in of the max is legal}

In general $\E[\max_i X_i] \ge \max_i \E[X_i]$---strict inequality.  So why is the push-in legal here?  Because $\pi_{k+1:N}$ may \emph{depend on $\bm{y}_k$} (through $b_{k+1}$): a policy that observes $\bm{y}_k$ before committing the continuation is strictly more powerful than one that commits in advance, and the problem allows the former.  Formally: maximizing pointwise inside the expectation gives the same value as maximizing over all measurable continuation plans outside.  This is the \emph{interchangeability lemma} for adaptive optimization---the rigorous form of Bellman's \emph{principle of optimality}.

\subsection{The Bellman system}

\begin{align}
  W_{N+1}(b_{N+1}; \bm{c}) &\;=\; H(p_0) - H(b_{N+1}),\\
  W_k(b_k; \bm{c}) &\;=\; \max_{\bm{s}_k}\; \E_{\bm{y}_k \mid b_k, \bm{s}_k, \bm{c}}\!\bigl[W_{k+1}(b_{k+1}; \bm{c})\bigr], \quad k = N, \ldots, 1,\\
  \pi^\star_k(b_k; \bm{c}) &\;=\; \argmax_{\bm{s}_k}\; \E_{\bm{y}_k \mid b_k, \bm{s}_k, \bm{c}}\!\bigl[W_{k+1}(b_{k+1}; \bm{c})\bigr],\\
  V_{\mathrm{adaptive}}(\bm{c}) &\;=\; W_1(p_0; \bm{c}).
\end{align}

\subsection{Summary}

Three ingredients yield the Bellman equation:
\begin{enumerate}[nosep, label=(\alph*)]
  \item Sufficient statistic reduction: replace $\bm{y}_{1:k-1}$ with $b_k$.
  \item Peeling: $\max_{\pi_{k:N}} = \max_{\bm{s}_k}\,\max_{\pi_{k+1:N}}$.
  \item Tower property with interchangeability: push $\max_{\pi_{k+1:N}}$ inside $\E_{\bm{y}_k}$.
\end{enumerate}
The recursion rests on two facts: the Markov property of beliefs under Bayesian updating, and the principle of optimality (the optimal continuation can depend on intervening observations without loss).

\end{document}
